\clearpage\pagestyle{otMain}\fontsize{11}{15}\selectfont
\corrigesChapitre{8}{Vecteurs du plan}
\begin{blocLivre}{cor-8-1}
\begin{corrige}{1}{Calculer un déplacement}
\demarche{Effectuer arrivée moins départ dans le bon ordre.}
On soustrait les coordonnées de l’origine à celles de l’extrémité : $\vect{AB}=(5-1;-1-2)=\res{(4;-3)}$. Dans l’autre sens, $\vect{BA}=\res{(-4;3)}=-\vect{AB}$. Les deux flèches ont la même longueur mais des sens opposés.
\label{fig-cor-8-1}
\figureGeo[corrige]{vect-01}
\end{corrige}
\end{blocLivre}
\begin{blocLivre}{cor-8-2}
\begin{corrige}{2}{Additionner des vecteurs}
\demarche{Additionner les déplacements composante par composante.}
Les coordonnées s’additionnent séparément : $\vec u+\vec v=(3-1;1+2)=\res{(2;3)}$. Sur le schéma, partir de l’origine, suivre $\vec u$, puis reporter $\vec v$ sans le tourner. L’arrivée est le point $(2;3)$ ; la flèche directe représente la somme.
\label{fig-cor-8-2}
\figureGeo[corrige]{vect-02}
\end{corrige}
\end{blocLivre}
\begin{blocLivre}{cor-8-3}
\begin{corrige}{3}{Multiplier par un réel}
\demarche{Interpréter séparément le signe du coefficient et sa valeur absolue.}
On multiplie chaque coordonnée : $3\vec u=\res{(6;-3)}$, de même sens que $\vec u$ et trois fois plus long. $-2\vec u=\res{(-4;2)}$ a le sens opposé et une norme double. Dans les deux cas, les vecteurs restent colinéaires à $\vec u$.
\label{fig-cor-8-3}
\figureGeo[corrige]{vect-03}
\end{corrige}
\end{blocLivre}
\begin{blocLivre}{cor-8-4}
\begin{corrige}{4}{Construire une image par translation}
\demarche{Ajouter le déplacement aux coordonnées du point de départ.}
La translation impose $\vect{AB}=(4;2)$. Il faut donc $x_B=-2+4$ et $y_B=1+2$, soit $\res{B(2;3)}$. Vérification : $\vect{AB}=(2-(-2);3-1)=(4;2)$. Ajouter le vecteur ne consiste pas à multiplier les coordonnées du point.
\label{fig-cor-8-4}
\figureGeo[corrige]{vect-04}
\end{corrige}
\end{blocLivre}
\begin{blocLivre}{cor-8-5}
\begin{corrige}{5}{Une norme}
\demarche{Utiliser la formule de norme dans un repère orthonormé.}
Par Pythagore, $\|\vec u\|=\sqrt{(-3)^2+4^2}=\sqrt{25}=\res5$. Le vecteur double vaut $(-6;8)$ et sa norme vaut $\sqrt{36+64}=\res{10}$. On retrouve $\|2\vec u\|=2\|\vec u\|$. Les coordonnées peuvent être négatives, la norme ne l’est pas.
\label{fig-cor-8-5}
\figureGeo[corrige]{vect-05}
\end{corrige}
\end{blocLivre}
\begin{blocLivre}{cor-8-6}
\begin{corrige}{6}{La relation de Chasles}
\demarche{Réorganiser la somme pour suivre les déplacements.}
Les déplacements successifs de $A$ à $B$, de $B$ à $C$ puis de $C$ à $D$ se regroupent en $\res{\vect{AD}}$. Dans la deuxième somme, on peut permuter les termes : $\vect{CA}+\vect{AB}=\res{\vect{CB}}$. L’ordre d’écriture des lettres doit correspondre à un chemin continu.
\label{fig-cor-8-6}
\figureGeo[corrige]{vect-06}
\end{corrige}
\end{blocLivre}
\begin{blocLivre}{cor-8-7}
\begin{corrige}{7}{Tester la colinéarité}
\demarche{Calculer un déterminant exact plutôt qu’apprécier la pente à l’œil.}
Le premier déterminant vaut $2\times9-3\times6=0$ ; les vecteurs sont colinéaires, et le second est trois fois le premier. Le second déterminant vaut $2\times7-3\times4=\res2\ne0$ ; ces deux vecteurs ne sont pas colinéaires. Des coordonnées « presque proportionnelles » ne suffisent pas.
\label{fig-cor-8-7}
\figureGeo[corrige]{vect-07}
\end{corrige}
\end{blocLivre}
\begin{blocLivre}{cor-8-8}
\begin{corrige}{8}{Trois points alignés ?}
\demarche{Comparer les vecteurs de même origine, puis leurs longueurs orientées.}
$\vect{AB}=(2;3)$ et $\vect{AC}=(4;6)=2\vect{AB}$ : les vecteurs sont colinéaires, donc les trois points sont alignés. De plus $\vect{BC}=(2;3)=\vect{AB}$, ce qui montre que $\res{B\text{ est le milieu de }[AC]}$. Le coefficient 2 donne ici davantage que le simple alignement.
\label{fig-cor-8-8}
\figureGeo[corrige]{vect-08}
\end{corrige}
\end{blocLivre}
\begin{blocLivre}{cor-8-9}
\begin{corrige}{9}{Le quatrième sommet}
\demarche{Choisir l’égalité de côtés opposés correspondant à l’ordre des sommets.}
Dans un parallélogramme, $\vect{AD}=\vect{BC}$. Or $\vect{BC}=(1;3)$ ; donc $D=A+(1;3)=\res{(1;4)}$. On peut contrôler avec $\vect{DC}=(4;1)=\vect{AB}$. L’ordre $A,B,C,D$ est nécessaire pour choisir la bonne égalité vectorielle.
\label{fig-cor-8-9}
\figureGeo[corrige]{vect-09}
\end{corrige}
\end{blocLivre}
\begin{blocLivre}{cor-8-10}
\begin{corrige}{10}{Une combinaison}
\demarche{Calculer d’abord les multiples, puis effectuer la soustraction.}
On calcule $2\vec u-\vec v=(4;2)-(-1;3)=\res{(5;-1)}$. L’expression cherchée est donc précisément $\res{2\vec u-\vec v}$. La soustraction d’un vecteur signifie ajouter son opposé ; les deux signes de ses coordonnées changent.
\label{fig-cor-8-10}
\figureGeo[corrige]{vect-10}
\end{corrige}
\end{blocLivre}
\begin{blocLivre}{cor-8-11}
\begin{corrige}{11}{Un paramètre de colinéarité}
\demarche{Traduire la condition géométrique en une équation factorisable.}
Le déterminant est $m(m-1)-6=m^2-m-6=(m-3)(m+2)$. Il est nul exactement pour $\res{m=3\text{ ou }m=-2}$. Le premier vecteur n’est jamais nul puisque sa deuxième coordonnée est 2 ; le critère s’applique sans ambiguïté. On contrôle les deux couples obtenus par proportionnalité.
\label{fig-cor-8-11}
\figureGeo[corrige]{vect-11}
\end{corrige}
\end{blocLivre}
\begin{blocLivre}{cor-8-12}
\begin{corrige}{12}{Un point sur une droite}
\demarche{Utiliser le déterminant sans imposer à tort l’appartenance au segment.}
$\vect{AB}=(4;2)$ et $\vect{AM}=(t-1;3)$. Le déterminant vaut $4\times3-2(t-1)=14-2t$. L’alignement équivaut à $14-2t=0$, d’où $\res{t=7}$. Le point $M(7;5)$ est sur le prolongement de $[AB]$, ce qui est permis : alignement ne signifie pas appartenance au segment.
\label{fig-cor-8-12}
\figureGeo[corrige]{vect-12}
\end{corrige}
\end{blocLivre}
\begin{blocLivre}{cor-8-13}
\begin{corrige}{13}{Une base non usuelle}
\demarche{Traduire une combinaison vectorielle par deux équations de coordonnées.}
Le déterminant vaut $-1-1=-2\ne0$. Pour la décomposition, il faut $a+b=4$ et $a-b=2$. En ajoutant les équations, $2a=6$ ; ainsi $a=3$ et $b=1$. Donc $\res{\vec w=3\vec u+\vec v}$. La non-colinéarité garantit l’unicité de cette décomposition.
\label{fig-cor-8-13}
\figureGeo[corrige]{vect-13}
\end{corrige}
\end{blocLivre}
\begin{blocLivre}{cor-8-14}
\begin{corrige}{14}{Les milieux d’un triangle}
\demarche{Comparer deux vecteurs directeurs des droites visées.}
$\vect{IJ}=(-2;2)$ et $\vect{BC}=(-4;4)$. On a $\vect{IJ}=\tfrac12\vect{BC}$ ; les vecteurs sont non nuls et colinéaires. Les droites $(IJ)$ et $(BC)$ sont donc parallèles. La même égalité donne $IJ=BC/2$, ce qui retrouve le théorème des milieux.
\label{fig-cor-8-14}
\figureGeo[corrige]{vect-14}
\end{corrige}
\end{blocLivre}
\begin{blocLivre}{cor-8-15}
\begin{corrige}{15}{Un déplacement composé}
\demarche{Additionner les déplacements puis tester simultanément les deux coordonnées.}
Le déplacement total est $n(2;1)+(n-1)(-1;2)=(n+1;3n-2)$. Après ajout du point de départ, $\res{M_n(n+2;3n-1)}$. Pour atteindre $(8;8)$, l’abscisse impose $n=6$, mais l’ordonnée vaudrait 17. Il n’existe donc aucun entier $n$ convenable. Les deux coordonnées doivent correspondre au même $n$.
\label{fig-cor-8-15}
\figureGeo[corrige]{vect-15}
\end{corrige}
\end{blocLivre}
